\subsection{Notational conventions.}
We adopt the standard Bachman-Landau, Hardy, and Vinogradov notations:
$f=O(g)$ and $f\ll g$ mean that there is a positive constant $C$
so that $|f| \le Cg$ throughout the domain of $f$.  The constant $C$ is independent of any parameters, unless specified by subscripts, e.g.
$f(x)=O_{\eps} (x^\eps)$.
Also, $f(x) \sim g(x)$ as $x\to \infty$ means $\lim_{x\to \infty} f(x)/g(x)=1$
and $f(x)=o(g(x))$ means that $\lim_{x\to\infty} f(x)/g(x)=0$.

For $\sigma\in \cS_n$, the notation
$\beta|\sigma$ means that $\beta$ is a \emph{divisor} of the permutation
$\sigma$, i.e. a product of some subset of the cycles of $\sigma$.
$|\beta|$ is the size (length) of $\beta$.

$\one(S)$ is the indicator function of statement $S$; $\displaystyle \one(S) = 
\begin{cases} 1 & S \text{ is true} \\ 0 & S \text{ is false}. \end{cases}$

